%------------------------------------------------------------------------------%
\begin{question}
  Encontre o maior e menor valores que a função
  \[
    f(x, y) = xy
  \]
  assume na elipse
  \[
    \frac{x^2}{8} + \frac{y^2}{2} = 1
  \]
\end{question}

%------------------------------------------------------------------------------%
\begin{resolution}{Interpretando o Problema}
  Queremos encontrar os valores extremos da função
  \[
    f(x, y) = xy
  \]
  \pause
  sujeito a restrição
  \[
    g(x, y) = \frac{x^2}{8} + \frac{y^2}{2} - 1 = 0
  \]

  \pause
  Para isso precisamos encontrar os valores $x$, $y$ e $\lambda$ tais que
  \[
    \nabla f = \lambda \nabla g
    \qquad\qquad
    g(x, y) = 0
  \]
\end{resolution}

%------------------------------------------------------------------------------%
\begin{resolution}{Gradientes}
\begin{columns}
\column{0.4\linewidth}
  Gradiente de $f$
  \[
    \pause
    \D{f}{x}
    \pause
    = \D{}{x}\left(xy\right)
    \pause
    = y
  \]
  \[
    \pause
    \D{f}{y}
    \pause
    = \D{}{y}\left(xy\right)
    \pause
    = x
  \]
  \pause
  \[
    {\color<19>{blue}\nabla f = \vetor{y}{x}}
  \]
\column{0.6\linewidth}
  \pause
  Gradiente de $g$
  \[
    \pause
    \D{g}{x}
    \pause
    = \D{}{x}\left(\frac{x^2}{8} + \frac{y^2}{2} - 1\right)
    \pause
    = \frac{2x}{8}
    \pause
    = \frac{x}{4}
  \]
  \[
    \pause
    \D{g}{y}
    \pause
    = \D{}{y}\left(\frac{x^2}{8} + \frac{y^2}{2} - 1 \right)
    \pause
    = \frac{2y}{2}
    \pause
    = y
  \]
  \pause
  \[
    {\color<19>{blue}\nabla g = \vetor{\dfrac{x}{4}}{y}}
  \]
\end{columns}
\end{resolution}

%------------------------------------------------------------------------------%
\begin{resolution}{Sistema}
\begin{columns}[t]
\column{0.4\linewidth}
  A equação
  \[
    \nabla f = \lambda \nabla g
  \]
  \pause
  se torna
  \[
    \vetor{y}{x} = \lambda \Vetor{\dfrac{x}{4}}{y}
  \]
  {\color<8>{blue}%
  \pause
  \[
    \begin{cases}
      4y = \lambda x \\[2mm]
      x = \lambda y
    \end{cases}
  \]}
\column{0.4\linewidth}
  \pause
  A equação
  \[
    g(x, y) = 0
  \]
  \pause
  se torna
  \[
    \frac{x^2}{8} + \frac{y^2}{2} - 1 = 0
  \]
  \[
    \pause
    \frac{x^2}{8} + \frac{y^2}{2} = 1
  \]
  {\color<8>{blue}%
  \pause
  \[
    x^2 + 4y^2 = 8
  \]}
\end{columns}
\end{resolution}

%------------------------------------------------------------------------------%
\begin{resolution}{Resolvendo o Sistema Não-Linear}
\begin{columns}
\column{0.5\linewidth}
  Precisamos resolver o sistema
  \[
    \begin{cases}
      4y = \lambda x \\
      x  = \lambda y \\
      x^2 + 4y^2 = 8
    \end{cases}
  \]
  \pause
  Substituindo \(x = \lambda y\) em \( 4y = \lambda x\)
  \[
    \pause
    4y = \lambda x
    \pause
    = \lambda \lambda y
  \]
  \[
    \pause
    4y = \lambda^2 y
    \pause
    \qquad
    \text{\color{red}Não divida por $y$!}
  \]
\column{0.5\linewidth}
  \pause
  Assim \(y=0\)
  \pause
  ou \(y\neq 0\)
  \pause
  e
  \begin{align*}
    \onslide<+->{\lambda^2     & = 4  \\}
    \onslide<+->{\abs{\lambda} & = 2  \\}
    \onslide<+->{\lambda       & = \pm 2}
  \end{align*}

  \onslide<+->{Caso 1: \emph{\(y=0\)}} \\
  \onslide<+->{Caso 2: \emph{\(y\neq 0\)} e \emph{\(\lambda=2\)}} \\
  \onslide<+->{Caso 3: \emph{\(y\neq 0\)} e \emph{\(\lambda=-2\)}}
\end{columns}
\end{resolution}

%------------------------------------------------------------------------------%
\begin{resolution}{Caso 1}
  \[
    y = 0
  \]
  \pause
  \[
    x = \lambda y
    \pause
    = 0
  \]
  \pause
  \[
    x^2 + 4y^2
    \pause
    = 0
    \pause
    \neq 8
  \]
  \pause
  Não resolve o sistema

  \pause
  \vspace{\baselineskip}
  Portanto \(y\neq 0\)
\end{resolution}

%------------------------------------------------------------------------------%
\begin{resolution}{Casos 2 e 3}
\begin{columns}
\column{0.5\linewidth}
  \[
    y \neq 0
    \qquad
    \lambda = 2
  \]
  \[
    \pause
    x
    = \lambda y
    \pause
    = 2 y
  \]
  \[
    \pause
    x^2 = 4y^2
  \]
  \pause
  \[
    x^2 + 4y^2
    \pause
    = 4y^2 + 4y^2
    \pause
    = 8y^2
    \pause
    = 8
  \]
  \[
    \pause
    y^2 = 1
  \]
  \pause
  Portanto \(y=\pm 1\)
  \pause
\column{0.5\linewidth}
  \[
    y \neq 0
    \qquad
    \lambda = -2
  \]
  \[
    \pause
    x
    = \lambda y
    \pause
    = -2 y
  \]
  \[
    \pause
    x^2 = 4y^2
  \]
  \pause
  \[
    x^2 + 4y^2
    \pause
    = 4y^2 + 4y^2
    \pause
    = 8y^2
    \pause
    = 8
  \]
  \[
    \pause
    y^2 = 1
  \]
  \pause
  Portanto \(y=\pm 1\)
\end{columns}
  \pause
  \vspace{\baselineskip}
  Temos então os pontos
  \(( 2,  1)\),
  \(( 2, -1)\),
  \((-2,  1)\) e
  \((-2, -1)\)
\end{resolution}

%------------------------------------------------------------------------------%
\begin{resolution}{Solução}
  \onslide<+->{Avaliando a função \(f(x, y)=xy\) nos pontos}
  \begin{align*}
    \onslide<+->{f(-2, -1)}
    \onslide<+->{& = (-2)(-1)}
    \onslide<+->{  =  2 \\}
    \onslide<+->{f(\hphantom{-}2, -1)}
    \onslide<+->{& = 2 (-1)}
    \onslide<+->{  = -2 \\}
    \onslide<+->{f(-2, \hphantom{-}1)}
    \onslide<+->{& = -2 \times 1}
    \onslide<+->{  = -2 \\}
    \onslide<+->{f(\hphantom{-}2, \hphantom{-}1)}
    \onslide<+->{& =  2 \times 1}
    \onslide<+->{  = 2}
  \end{align*}

  \onslide<+->{O valor máximo de $f$ é $\hpm 2$
    e ocorre nos pontos \((-2, -1)\) e \(( 2, 1)\)}

  \onslide<+->{O valor mínimo de $f$ é $-2$
    e ocorre nos pontos \((-2, 1)\) e \((2, -1)\)}
\end{resolution}

%------------------------------------------------------------------------------%
